\section{La función índice. Versión general de la fórmula integral de Cauchy}
\subsection{La función índice}

Observemos que:
\[
\frac{1}{2\pi i} \int_{C(0,1)} \frac{1}{w - z}\,dw =
\begin{cases}
    \displaystyle \frac{1}{2\pi i} \int_{\partial D(0,1)} \frac{f(w)}{w - z}\,dw = f(z) = 1, & \text{por la fórmula integral de Cauchy, si } |z| < 1, \\[1em]
    0, & \text{por el teorema de Cauchy, si } |z| > 1.
\end{cases}
\]

En efecto, cuando $|z| < 1$ podemos aplicar la fórmula integral de Cauchy con $f(w) = 1$, obteniendo $f(z) = 1$.  
Por otro lado, si $|z| > 1$, la función $w \mapsto \dfrac{1}{w - z}$ es holomorfa en $\mathbb{C} \setminus \{z\}$, y la curva $C(0,1)$ es homótopa a cero en dicho dominio. Por el teorema de Cauchy, la integral se anula.

Este razonamiento motiva la siguiente definición.

\begin{definición} [Función índice]
Sea $\gamma : [a,b] \to \mathbb{C}$ una curva cerrada en $\mathbb{C}$.  
Definimos la \textbf{función índice de $\gamma$} como:
\[
\Ind_\gamma(z) := \frac{1}{2\pi i} \int_\gamma \frac{1}{w - z}\,dw, 
\quad \forall z \in \mathbb{C} \setminus \gamma^*.
\]
\end{definición}

Observemos que si $\gamma$ es una curva cerrada en $\mathbb{C}$, entonces $\mathbb{C} \setminus \gamma^*$ es un abierto que tendrá una o varias componentes conexas, y solo una de ellas no acotada.\\

Como $\gamma^*$ es compacto, existe $R > 0$ tal que $\gamma^* \subset D(0,R)$, luego $\mathbb{C} \setminus \overline{D}(0,R) \subset \mathbb{C} \setminus \gamma^*$, y como $\mathbb{C} \setminus \overline{D}(0,R)$ es conexo, entonces $\mathbb{C} \setminus \overline{D}(0,R)$ está contenido en una componente conexa de $\mathbb{C} \setminus \gamma^*$, que será la única componente conexa no acotada de $\mathbb{C} \setminus \gamma^*$, ya que el resto estarán  contenidas en $\overline{D}(0,R)$, que es acotado.\\

\begin{proposición}[Propiedades de la función índice]
    Sea $\gamma : [a,b] \to \mathbb{C}$ una curva cerrada en $\mathbb{C}$. Entonces:
    \begin{enumerate}
        \item $\Ind_\gamma$ es continua en $\mathbb{C} \setminus \gamma^*$.
        \item $\Ind_\gamma$ solo toma valores enteros, es decir, $\Ind_\gamma (\mathbb{C} \setminus \gamma^*) \subset \mathbb{Z}$.
        \item $\Ind_\gamma$ es constante en cada componente conexa de $\mathbb{C} \setminus \gamma^*$.
        \item $\Ind_\gamma \equiv 0$ en la componente conexa no acotada de $\mathbb{C} \setminus \gamma^*$.
    \end{enumerate}
\end{proposición}

\begin{proof}
    Veamos cada una de las propiedades:
    \begin{enumerate}
        \item Sea $z_1 \in \mathbb{C} \setminus \gamma^*$. Entonces, como $\mathbb{C} \setminus \gamma^*$ es abierto, existe $\delta > 0$ tal que $D(z_1, \delta) \subset \mathbb{C} \setminus \gamma^*$.\\
        Aplicando el lema 2 del tema 8, existen $t_0, t_1, \ldots, t_n \in [a,b]$ particion que cumplen que $\forall j \in \{1, \ldots, n\}$ se tiene que $\gamma([t_{j-1}, t_j])$ está contenida en un disco abierto $D_j \subset \mathbb{C} \setminus \overline{D} (z_1, \delta)$.\\
        Por esto mismo, podemos obtener una curva $C^1$ o $C^1$ a trozos llamada $\Gamma$ mediante un numero finito de deformaciones elementales de $\gamma$ en $\mathbb{C} \setminus \overline{D}(z_1, \delta)$, pues para cada $z \in D(z_1, \delta)$, $\gamma$ la función $w \mapsto \frac{1}{w - z}$ es holomorfa en $\mathbb{C} \setminus \overline{D}(z_1, \delta)$.\\
        Es claro que $\int_\gamma \frac{1}{w - z} dw = \int_\Gamma \frac{1}{w - z} dw \, \quad \forall z \in D(z_1, \delta)$.\\
        Ahora, si $z \in D(z_1, \delta/2)$, entonces:
        \[
        \left| \Ind_\gamma (z) - \Ind_\gamma (z_1) \right| = 
        \left| \left( \frac{1}{2 \pi i} \int_\gamma \frac{dw}{w - z} \right) - 
        \left( \frac{1}{2 \pi i} \int_\gamma \frac{dw}{w - z_1} \right) \right| =
        \frac{1}{2 \pi i} \left| \int_\Gamma \frac{1}{w - z} - \int_\Gamma \frac{1}{w - z_1} dw \right| =  
        \]
        \[
        = \frac{1}{2 \pi } \left| \int_\Gamma \frac{(w - z_1) - (w - z)}{(w - z)(w - z_1)} dw \right| =
        \frac{1}{2 \pi } \left| \int_\Gamma \frac{z - z_1}{(w - z)(w - z_1)} dw \right| \leq
        \]
        \[
        \frac{1}{2 \pi } \Long (\Gamma) \left(\sup_{w \in \Gamma} \frac{|z-z_1|}{|w - z||w - z_1|}\right)  =
        \frac{1}{2 \pi } \Long (\Gamma) |z - z_1| \left(\sup_{w \in \Gamma} \frac{1}{|w - z||w - z_1|}\right)  \leq^*
        \frac{1}{\pi  \delta^2} \Long (\Gamma) |z - z_1|
        \]
        Usando en $*$ que como $\Gamma \subset \mathbb{C} \setminus \overline{D}(z_1, \delta)$, entonces $\forall w \in \Gamma$ se tiene que $|w - z_1| \geq \delta$, y como $|w-z| \geq |w - z_1| - |z - z_1| \geq \delta - \delta/2 = \delta/2$, entonces se cumple que $\frac{1}{|w - z|} \leq \frac{2}{\delta}$. Luego 
        \[
            \sup_{w \in \Gamma} \frac{1}{|w-z||w-z_1|} \leq \frac{1}{(\delta/2)\delta} = \frac{2}{\delta^2}
        \]
        Así, finalmente, tenemos que:
        \[
        \left| \Ind_\gamma(z) - \Ind_\gamma(z_1) \right| \leq
        \frac{\Long (\Gamma)}{\pi \delta^2} |z - z_1| \xrightarrow{z \to z_1} 0 \quad \forall z \in D(z_1, \delta/2)
        \]
        \item Sabemos que $\Ind_\gamma (z) = \frac{1}{2 \pi i } \int_\gamma \frac{1}{w - z} dw \, \quad \forall z \in \mathbb{C} \setminus \gamma^*$.\\
        Ahora sea $z \in \mathbb{C} \setminus \gamma^*$:
        \[
        \Ind_\gamma(z) = \frac{1}{2 \pi i} \int_\gamma \frac{1}{w - z} dw =
        \frac{1}{2 \pi i} \int_a^b \frac{\gamma'(t)}{\gamma(t) - z} dt
        \]
        Definamos:
        \[
        \begin{array}{ccc}
            \varphi : [a,b] & \to & \mathbb{C} \\
            x & \mapsto & \int_a^x \frac{\gamma'(t)}{\gamma(t) - z} dt
        \end{array}
        \]
        Y tengamos:
        \[
        h(x) = e^{-\varphi(x)} (\gamma(x) - z) \, \quad \forall x \in [a,b]
        \]
        Entonces se cumple que la derivada de $h$ es:
        \[
        h'(x) = 
        e^{-\varphi(x)} \gamma'(x) - \varphi'(x) e^{-\varphi(x)} (\gamma(x) - z) =
        e^{-\varphi(x)} \left( \gamma'(x) - \underbrace{\frac{\gamma'(x)}{\gamma(x) - z}}_{\varphi'(x)} (\gamma(x) - z) \right) =
        0 \, \quad \forall x \in [a,b]
        \]
        Por lo que:
        \[
        h \text{ es constante en } [a,b] \implies 
        h(a) = h(b) \implies 
        e^{-\varphi(a)} (\gamma(a) - z) = e^{-\varphi(b)} (\gamma(b) - z) \implies_{\gamma \text{ cerrada}}
        \]
        \[
        e^0 (\gamma(a) - z) = 
        e^{-\varphi(b)} (\gamma(a) - z) \implies_{z \notin \gamma^* \implies \gamma(a) \neq z} e^{0} = e^{- \varphi(b)}
        \]
        Como la función exponencial es periódica con periodo $2 \pi i$, se tiene que:
        \[
        \exists k \in \mathbb{Z} : - \varphi (b) =
        2k\pi i \implies
        \varphi (b) = 2(-k) \pi i \implies
        \Ind_\gamma (z) = \frac{1}{2 \pi i} \varphi (b) = -k \in \mathbb{Z}
        \]
        \item Veamos que:
        \[
        \begin{cases}
            \Ind_\gamma \text{ es continua} \\
            \Omega \text{ componente conexa de } \mathbb{C} \setminus \gamma^*
        \end{cases} \implies
        \Ind_\gamma (\Omega) \text{ es conexo } \implies_{\Ind_\gamma (\Omega) \subset \mathbb{Z}}
        \]
        \[
        \Ind_\gamma \text{ es un unico valor de } \mathbb{Z} \implies
        \Ind_\gamma \text{ es constante en } \Omega
        \]
        \item Como $\gamma^*$ es compacto, existe $R > 0$ tal que $\gamma^* \subset D(0,R)$, sea ahora $R' > R$:
        \[
        \left| \Ind_\gamma (R) \right| \underbrace{=}_{R,2R' \in \text{ componente conex. no acotada de } \mathbb{C} \setminus \gamma^*}
        \left| \Ind_\gamma (2R') \right| =
        \left| \frac{1}{2 \pi i} \int_\gamma \frac{dw}{w - 2R'} \right| \leq
        \]
        \[
        \frac{1}{2 \pi} \Long (\gamma) \sup_{w \in \gamma^*} \frac{1}{|w - 2R'|} \leq
        \frac{1}{2 \pi} \Long (\gamma) \frac{1}{2R'} \xrightarrow{R' \to +\infty} 0
        \]
        Ya que $w \in \gamma \implies |w| \leq R < R' \implies |w - 2R'| \geq R' \implies \frac{1}{|w - 2R'|} \leq \frac{1}{R'}$.\\
    \end{enumerate}
\end{proof}

\begin{definición}
    Sea $\gamma : [a,b] \to \mathbb{C}$ una curva cerrada en $\mathbb{C}$. Se define el \textbf{interior} y \textbf{exterior} de $\gamma$ como:
    \[
        \Int(\gamma) := \{ z \in \mathbb{C} \setminus \gamma^* : \Ind_\gamma (z) \neq 0 \},
        \quad
        \Ext(\gamma) := \{ z \in \mathbb{C} \setminus \gamma^* : \Ind_\gamma (z) = 0 \}.
    \]
\end{definición}

\textbf{Intuición geométrica de la función índice:} Si $\gamma$ es una curva cerrada en $\mathbb{C}$ y $z_0 \in \mathbb{C} \setminus \gamma^*$, entonces $Ind_\gamma (z)$ mide el número de vueltas que da la curva $\gamma$ alrededor del punto $z$, teniendo en cuenta la orientación de la curva

Para todo $t \in [a,b]$, escribamos $w(t) = \gamma(t)$ y expresemos
\[
    w(t) - z_0 = r(t)e^{i\theta(t)},
\]
donde $r(t) > 0$ y $\theta(t)$ es una función continua que representa el argumento de $w(t) - z_0$.

Derivando respecto de $t$ se obtiene
\[
    \dot{w}(t) = e^{i\theta(t)}\big(r'(t) + i\,r(t)\theta'(t)\big),
\]
y por tanto
\[
    \frac{\dot{w}(t)}{w(t) - z_0} = \frac{r'(t)}{r(t)} + i\,\theta'(t).
\]

Integrando a lo largo de la curva $\gamma$ se tiene
\[
    \int_\gamma \frac{dw}{w - z_0}
    = \int_a^b \frac{r'(t)}{r(t)}\,dt + i \int_a^b \theta'(t)\,dt.
\]
El primer término es nulo, ya que
\[
    \int_a^b \frac{r'(t)}{r(t)}\,dt = \log r(b) - \log r(a) = 0,
\]
porque $\gamma$ es cerrada y, por tanto, $r(a)=r(b)$.

El segundo término corresponde al cambio total del ángulo $\theta(t)$ cuando la curva recorre una vuelta completa, es decir,
\[
    \int_a^b \theta'(t)\,dt = \Delta \theta = 2\pi n,
\]
donde $n \in \mathbb{Z}$ es el número de vueltas que $\gamma$ da alrededor de $z_0$, contado con signo.

Sustituyendo en la expresión anterior se obtiene
\[
    \int_\gamma \frac{dw}{w - z_0} = i(2\pi n),
\]
y por tanto
\[
    \Ind_\gamma(z_0) = \frac{1}{2\pi i} \int_\gamma \frac{dw}{w - z_0} = n.
\]
En consecuencia, el índice representa el número de vueltas (positivas o negativas) que la curva $\gamma$ da alrededor de $z_0$.


\begin{observación}
    Puede ocurrir que $Int(\gamma) = \emptyset$. Por ejemplo, si $\gamma$ es un segmento de ida y vuelta, ó si $\gamma = \Gamma \cup \Gamma^-$, siendo $\Gamma$ una curva simple con exteriores distintos.\\
    Si $\gamma$ es una curva cerrada simple, orientada positivamente entonces
    \[
        \Ind_\gamma (z) = \begin{cases}
            1, & \text{si } z \in \Int(\gamma), \\
            0, & \text{si } z \in \Ext(\gamma).
        \end{cases}
    \]
\end{observación}



